
\documentclass[final,1p,times]{elsarticle}

% elsarticle already loads geometry
\geometry{
    left=0.55in,
    right=0.55in,
    top=0.40in,
    bottom=0.45in
}

\usepackage{amsmath}
\usepackage{amssymb}
\usepackage{bm}
\usepackage{algorithm}
\usepackage{algpseudocode}
\allowdisplaybreaks

\newcommand{\ket}[1]{\left|#1\right\rangle}
\newcommand{\bra}[1]{\left\langle#1\right|}
\newcommand{\Tr}{\operatorname{Tr}}
\newcommand{\Var}{\operatorname{Var}}
\newcommand{\Cov}{\operatorname{Cov}}
\newcommand{\argmin}{\operatorname*{arg\,min}}

\algrenewcommand\algorithmicrequire{\textbf{Input:}}
\algrenewcommand\algorithmicensure{\textbf{Output:}}

% ---------------------------------------------------------------------
% Breakable, non-floating algorithm environment.
% Unlike the standard algorithm float, this environment:
% (i) stays exactly where it is written in the source, and
% (ii) can continue across page boundaries.
% ---------------------------------------------------------------------
\makeatletter
\newenvironment{breakablealgorithm}
{%
  \par\addvspace{0.75\baselineskip}%
  \refstepcounter{algorithm}%
  \hrule height 0.6pt depth 0pt
  \kern 3pt
  \renewcommand{\caption}[2][\relax]{%
    \noindent\textbf{Algorithm~\thealgorithm. ##2}\par
    \ifx\relax##1\relax
      \addcontentsline{loa}{algorithm}{\protect\numberline{\thealgorithm}##2}%
    \else
      \addcontentsline{loa}{algorithm}{\protect\numberline{\thealgorithm}##1}%
    \fi
    \kern 3pt
    \hrule height 0.4pt depth 0pt
    \kern 4pt
  }%
}
{%
  \kern 4pt
  \hrule height 0.6pt depth 0pt
  \par\addvspace{0.75\baselineskip}%
}
\makeatother

\begin{document}

\begin{frontmatter}

\title{Critical Phase Shadow Reservoirs}

\author{Chinmoy Biswas}

\end{frontmatter}

% =====================================================================
\section{Methodological Overview}
\label{sec:overview}
% =====================================================================

The central object studied in this work is a hardware-aware,
window-encoded quantum reservoir implemented on an IBM Heron-compatible
coupling graph. The proposed Critical Phase Shadow Reservoir (PCSR) is
not a separate reservoir architecture introduced after the baseline
model; rather, it is a controlled extension of the same reservoir
backbone. The baseline uses full controlled-$Z$ interactions on the
selected topology, whereas PCSR replaces each controlled-$Z$ gate by a
tunable controlled-phase interaction $CP(2g)$. Thus, the interaction
parameter $g$ continuously connects a non-interacting reservoir
($g=0$) to the full-CZ limit ($g=\pi/2$).

The methodology is organized around one question: how does the
interaction regime of a hardware-constrained quantum reservoir affect
(i) information spreading and entanglement, (ii) temporal memory, and
(iii) the usefulness of the resulting quantum features for supervised
prediction? Classical-shadow-style observable extraction is then used
to connect the quantum reservoir to the LUQPI readout and to test
finite-shot and noisy-hardware robustness.

The complete logical flow is

\begin{equation}
\boxed{
\text{hardware topology}
\rightarrow
\text{window encoding}
\rightarrow
U_{\mathcal R}(g)
\rightarrow
\ket{\psi_t(g)}
\rightarrow
\text{Pauli/shadow features}
\rightarrow
\text{LUQPI readout}
}
\label{eq:main_pipeline}
\end{equation}

The following sections first define the common reservoir backbone and
benchmark tasks. PCSR is then introduced only as the critical-phase
generalization of that backbone. This ordering avoids redefining the
same temporal encoder, topology, and evaluation metrics in each
experiment.

% =====================================================================
\section{Shared Hardware-Aware Reservoir Backbone}
\label{sec:backbone}
% =====================================================================

\subsection{Hardware topology selection}

Let the device coupling graph be
$\mathcal G=(\mathcal V,\mathcal E_{\mathrm{HW}})$, where vertices are
physical qubits and edges are native two-qubit couplers. A reservoir
topology is represented by the selected physical-qubit set
$\mathcal Q$ and the corresponding interaction edge set
$\mathcal E_{\mathcal R}$. The implementation considers disconnected,
zigzag, star-hub, single-hub chain, double-hub ladder, and heavy-hex
ring motifs. Hardware-native constructions preferentially select
low-readout-error qubits while preserving the desired connectivity.
The star-hub configuration is retained as a comparison topology and
records the number of requested edges that are not native.

\begin{breakablealgorithm}
\caption{Hardware-aware reservoir topology construction}
\label{alg:topology}
\begin{algorithmic}[1]

\Require Hardware graph
$\mathcal{G}=(\mathcal{V},\mathcal{E}_{\mathrm{HW}})$,
reservoir size $N$, topology $\mathcal{T}$,
readout errors $\epsilon_q$

\Ensure Selected qubits $\mathcal{Q}$ and
reservoir edges $\mathcal{E}_{\mathcal{R}}$

\If{$\mathcal{T}=\mathrm{Disconnected}$}
    \State $\mathcal{Q}\gets$
    low-error native path of $N$ qubits
    \State $\mathcal{E}_{\mathcal{R}}\gets\emptyset$

\ElsIf{$\mathcal{T}=\mathrm{Zigzag}$}
    \State $\mathcal{Q}\gets$
    low-error native path of $N$ qubits
    \State $\mathcal{E}_{\mathcal{R}}
    \gets$ native edges induced by $\mathcal{Q}$

\ElsIf{$\mathcal{T}=\mathrm{StarHub}$}
    \State Select lowest-error degree-$3$ qubit $h$
    \State Expand $\mathcal{Q}$ from $h$ until
    $|\mathcal{Q}|=N$
    \State $\mathcal{E}_{\mathcal{R}}
    \gets \{(h,q):q\in\mathcal{Q}\setminus\{h\}\}$
    \State Count non-native edges as routing overhead

\ElsIf{$\mathcal{T}=\mathrm{Chain}$}
    \State Select lowest-error degree-$3$ hub $h$
    \State $\mathcal{Q}\gets\{h\}$

    \While{$|\mathcal{Q}|<N$}
        \State Construct unselected neighboring frontier
        $\mathcal{F}$
        \State Select $q^\star\in\mathcal{F}$ favoring
        high degree, connectivity, and low error
        \State $\mathcal{Q}\gets
        \mathcal{Q}\cup\{q^\star\}$
    \EndWhile

    \State $\mathcal{E}_{\mathcal{R}}
    \gets$ native edges induced by $\mathcal{Q}$

\ElsIf{$\mathcal{T}=\mathrm{Ladder}$}
    \State Find degree-$3$ hubs $(h_1,h_2)$ satisfying
    $d_{\mathcal{G}}(h_1,h_2)=2$
    \State Select pair minimizing
    $\epsilon_{h_1}+\epsilon_{h_2}$
    \State Find common bridge qubit $b$
    \State $\mathcal{Q}\gets\{h_1,b,h_2\}$

    \While{$|\mathcal{Q}|<N$}
        \State Add a neighboring qubit with lowest
        readout error
    \EndWhile

    \State $\mathcal{E}_{\mathcal{R}}
    \gets$ native edges induced by $\mathcal{Q}$

\ElsIf{$\mathcal{T}=\mathrm{HexRing}$}
    \State Find hardware cycles with length $\geq 12$
    \State Select
    $C^\star=
    \argmin_C\sum_{q\in C}\epsilon_q$
    \State $\mathcal{Q}\gets$ ordered qubits of $C^\star$
    \State $\mathcal{E}_{\mathcal{R}}
    \gets$ native edges induced by $\mathcal{Q}$

\EndIf

\State \Return
$(\mathcal{Q},\mathcal{E}_{\mathcal{R}})$

\end{algorithmic}
\end{breakablealgorithm}

\subsection{Temporal-window encoding}

All reservoir variants in this work use the same temporal encoder. For
time step $t$, the latest $W$ samples are right-aligned in the
zero-padded window

\begin{equation}
\bm w_t =
[u_{t-W+1},\ldots,u_t],
\label{eq:shared_window}
\end{equation}

with unavailable early-time samples replaced by zero. Window position
$w$ is mapped cyclically to reservoir qubit

\begin{equation}
q(w)=w\bmod N.
\label{eq:shared_slot}
\end{equation}

The position weights $\phi_w$ are linearly spaced over $[0.5,1]$. The
accumulated rotation angle on qubit $i$ is

\begin{equation}
\theta_i^{(t)}
=
\pi
\sum_{\substack{w=0\\q(w)=i}}^{W-1}
w_{t,w}\phi_w ,
\label{eq:shared_theta}
\end{equation}

and the input-dependent encoder is

\begin{equation}
U_{\mathrm{enc}}(\bm w_t)
=
\bigotimes_{i=1}^{N}
R_y\!\left(\theta_i^{(t)}\right).
\label{eq:shared_encoder}
\end{equation}

This explicit temporal window is the source of history dependence in
the present implementation. The reservoir is initialized from
$\ket{0}^{\otimes N}$ for every $t$ rather than propagating
$\ket{\psi_{t-1}}$ into the next step.

\subsection{Baseline full-CZ reservoir}

The baseline reservoir applies the native topology-dependent CZ layer
followed by fixed local rotations. With fixed biases
$\beta_i^{(z)}$ and $\beta_i^{(x)}$, one step can be written as

\begin{equation}
U_{\mathrm{base}}
=
\left[
\prod_{i=1}^{N}
R_x(\beta_i^{(x)})
R_z(\beta_i^{(z)})
\right]
\left[
\prod_{(i,j)\in\mathcal E_{\mathcal R}}
CZ_{ij}
\right].
\label{eq:baseline_unitary}
\end{equation}

For $R$ repetitions,

\begin{equation}
U_{\mathcal R}^{\mathrm{base}}
=
U_{\mathrm{base}}^R ,
\end{equation}

and the state generated from the common encoder is

\begin{equation}
\ket{\psi_t^{\mathrm{base}}}
=
U_{\mathcal R}^{\mathrm{base}}
U_{\mathrm{enc}}(\bm w_t)
\ket{0}^{\otimes N}.
\label{eq:baseline_state}
\end{equation}

Algorithm~\ref{alg:qrc} gives the implementation of this baseline
window-encoded reservoir.

\begin{breakablealgorithm}
\caption{Sliding-window quantum reservoir evolution}
\label{alg:qrc}
\begin{algorithmic}[1]

\Require Input sequence
$\bm{u}=\{u_1,\ldots,u_T\}$,
topology $(\mathcal{Q},\mathcal{E}_{\mathcal{R}})$,
number of qubits $N$, window size $W$,
reservoir repetitions $R$

\Ensure Reservoir feature matrix $\bm{X}$

\State Draw fixed random biases
$\beta_i^{(z)}\sim\mathcal{U}(0,2\pi)$ and
$\beta_i^{(x)}\sim\mathcal{U}(0.4,1.2)$

\State Construct topology-dependent CZ circuit
from $\mathcal{E}_{\mathcal{R}}$

\For{$i=1,\ldots,N$}
    \State Apply $R_z(\beta_i^{(z)})$
    followed by $R_x(\beta_i^{(x)})$
\EndFor

\State Obtain fixed reservoir unitary
$U_{\mathrm{fixed}}$

\State
$U_{\mathcal{R}}\gets
(U_{\mathrm{fixed}})^R$

\State Define
$q(w)=w\bmod N$

\State Define temporal weights
$\phi_w\in[0.5,1]$

\For{$t=1,\ldots,T$}

    \State Construct zero-padded window
    $\bm{w}_t$ containing the latest $W$ inputs

    \State $\theta_i\gets0$,
    $i=1,\ldots,N$

    \For{$w=0,\ldots,W-1$}
        \State $i\gets q(w)$
        \State
        $\theta_i\gets
        \theta_i+\pi w_{t,w}\phi_w$
    \EndFor

    \State Construct
    \[
    U_{\mathrm{enc}}(\bm{w}_t)
    =
    \bigotimes_{i=1}^{N}R_y(\theta_i)
    \]

    \State
    \[
    \ket{\psi_t}
    \gets
    U_{\mathcal{R}}
    U_{\mathrm{enc}}(\bm{w}_t)
    \ket{0}^{\otimes N}
    \]

    \If{density-matrix simulation}
        \State
        $\rho_t\gets
        \ket{\psi_t}\bra{\psi_t}$
        \State Apply selected noise channels
        \State
        $\bm{x}_t\gets
        \textsc{ExtractFeatures}(\rho_t)$
    \Else
        \State
        $\bm{x}_t\gets
        \textsc{ExtractFeatures}(\ket{\psi_t})$
    \EndIf

    \State
    $\bm{X}_{t,:}\gets\bm{x}_t$

\EndFor

\State \Return $\bm{X}$

\end{algorithmic}
\end{breakablealgorithm}

\subsection{Baseline observable representation and readout}

For the baseline comparison, each state is represented using all local
Pauli expectations and pairwise $ZZ$ correlations,

\begin{equation}
\bm x_t^{\mathrm{base}}
=
[
\langle X_1\rangle,\langle Y_1\rangle,\langle Z_1\rangle,
\ldots,
\langle X_N\rangle,\langle Y_N\rangle,\langle Z_N\rangle,
\{\langle Z_iZ_j\rangle\}_{i<j}
],
\label{eq:baseline_features}
\end{equation}

with dimension

\begin{equation}
D_{\mathrm{base}}
=
3N+\binom{N}{2}.
\end{equation}

\begin{breakablealgorithm}
\caption{Quantum reservoir feature extraction}
\label{alg:feature}
\begin{algorithmic}[1]

\Require Quantum state $\ket{\psi}$ or
density matrix $\rho$, number of qubits $N$

\Ensure Feature vector
$\bm{x}\in
\mathbb{R}^{3N+N(N-1)/2}$

\State $\bm{x}\gets[\;]$

\For{$i=1,\ldots,N$}

    \If{state-vector representation}
        \State
        $x_i\gets\bra{\psi}X_i\ket{\psi}$
        \State
        $y_i\gets\bra{\psi}Y_i\ket{\psi}$
        \State
        $z_i\gets\bra{\psi}Z_i\ket{\psi}$
    \Else
        \State
        $x_i\gets\Tr(\rho X_i)$
        \State
        $y_i\gets\Tr(\rho Y_i)$
        \State
        $z_i\gets\Tr(\rho Z_i)$
    \EndIf

    \State Append
    $(x_i,y_i,z_i)$ to $\bm{x}$

\EndFor

\For{$i=1,\ldots,N$}
    \For{$j=i+1,\ldots,N$}

        \If{state-vector representation}
            \State
            $c_{ij}\gets
            \bra{\psi}Z_iZ_j\ket{\psi}$
        \Else
            \State
            $c_{ij}\gets
            \Tr(\rho Z_iZ_j)$
        \EndIf

        \State Append $c_{ij}$ to $\bm{x}$

    \EndFor
\EndFor

\State \Return $\bm{x}$

\end{algorithmic}
\end{breakablealgorithm}

Only the classical readout is trained. Ridge regression is used where
a linear readout is required,

\begin{equation}
\hat{\bm w}
=
\argmin_{\bm w}
\left[
\|\bm X_{\mathrm{tr}}\bm w-\bm y_{\mathrm{tr}}\|_2^2
+
\alpha\|\bm w\|_2^2
\right],
\label{eq:ridge}
\end{equation}

and predictive performance is reported using

\begin{equation}
\mathrm{NRMSE}
=
\sqrt{
\frac{
\operatorname{mean}
[(\hat{\bm y}-\bm y)^2]
}{
\Var(\bm y)+10^{-12}
}
}.
\label{eq:nrmse}
\end{equation}

Two reservoir-level diagnostics are used repeatedly later. The
delay-$k$ memory coefficient measures how accurately the current
feature vector reconstructs $u_{t-k}$,

\begin{equation}
MC_k
=
\frac{
\Cov(\hat u_{t-k},u_{t-k})^2
}{
\Var(\hat u_{t-k})\Var(u_{t-k})+10^{-12}
},
\label{eq:mc}
\end{equation}

while kernel quality is the numerical rank of the stacked state or
feature matrix, with singular values retained when
$\sigma_i>10^{-6}\sigma_1$. These definitions are used throughout the
PCSR experiments without being redefined in each subsection.

\begin{breakablealgorithm}
\caption{Classical readout and memory-capacity evaluation}
\label{alg:memory}
\begin{algorithmic}[1]

\Require Reservoir features $\bm{X}$,
target $\bm{y}$, input $\bm{u}$,
regularization $\alpha$, maximum delay $K$

\Ensure Prediction $\hat{\bm{y}}$,
NRMSE, and $\{MC_k\}_{k=1}^{K}$

\State Split $\bm{X}$ and $\bm{y}$
into training and testing sets

\State Train ridge regression:
\[
\hat{\bm{w}}
=
\argmin_{\bm{w}}
\left[
\|\bm{X}_{\mathrm{tr}}\bm{w}
-\bm{y}_{\mathrm{tr}}\|_2^2
+
\alpha\|\bm{w}\|_2^2
\right]
\]

\State
$\hat{\bm{y}}\gets
\bm{X}_{\mathrm{te}}\hat{\bm{w}}$

\State
\[
\mathrm{NRMSE}
\gets
\sqrt{
\frac{
\mathrm{mean}
[(\hat{\bm{y}}-\bm{y}_{\mathrm{te}})^2]
}{
\Var(\bm{y}_{\mathrm{te}})+10^{-12}
}
}
\]

\For{$k=1,\ldots,K$}

    \State
    $\bm{y}^{(k)}
    \gets
    \{u_1,\ldots,u_{T-k}\}$

    \State
    $\bm{X}^{(k)}
    \gets
    \{\bm{x}_{k+1},\ldots,\bm{x}_T\}$

    \State Align the two sequences

    \State Train ridge model
    $\bm{X}^{(k)}\rightarrow\bm{y}^{(k)}$

    \State Predict delayed input
    $\hat{\bm{u}}^{(k)}$

    \State
    \[
    MC_k
    \gets
    \frac{
    \Cov(\hat{\bm{u}}^{(k)},\bm{u}^{(k)})^2
    }{
    \Var(\hat{\bm{u}}^{(k)})
    \Var(\bm{u}^{(k)})
    +10^{-12}
    }
    \]

\EndFor

\State \Return
$\hat{\bm{y}}$, NRMSE,
$\{MC_k\}_{k=1}^{K}$

\end{algorithmic}
\end{breakablealgorithm}

\begin{breakablealgorithm}
\caption{Effective kernel-quality evaluation}
\label{alg:kernel}
\begin{algorithmic}[1]

\Require Reservoir representations
$\{\bm{s}_1,\ldots,\bm{s}_M\}$

\Ensure Effective kernel rank $K_Q$

\State Construct
\[
\bm{S}
\gets
[\bm{s}_1,\bm{s}_2,\ldots,\bm{s}_M]
\]

\State Compute singular values
\[
\{\sigma_i\}
\gets
\operatorname{SVD}(\bm{S})
\]

\State
$\sigma_{\max}\gets\max_i\sigma_i$

\State $K_Q\gets0$

\For{each $\sigma_i$}
    \If{$\sigma_i>10^{-6}\sigma_{\max}$}
        \State $K_Q\gets K_Q+1$
    \EndIf
\EndFor

\State \Return $K_Q$

\end{algorithmic}
\end{breakablealgorithm}
% =====================================================================
\section{Benchmark Tasks}
\label{sec:benchmark_tasks}
% =====================================================================

The same benchmark family is used for the baseline reservoir and for
PCSR, allowing changes in performance to be attributed to the
reservoir/readout construction rather than to different data. Unless
otherwise stated, the input is reproducibly generated as

\begin{equation}
u_t\sim\mathcal U(0,1),\qquad t=0,\ldots,T-1.
\label{eq:random_input}
\end{equation}

\begin{breakablealgorithm}
\caption{Random input sequence generation}
\label{alg:random_input}
\begin{algorithmic}[1]

\Require Sequence length $T$, random seed $s$
\Ensure Input sequence
$\bm{u}=\{u_0,u_1,\ldots,u_{T-1}\}$

\State Initialize pseudorandom generator using seed $s$

\For{$t=0,\ldots,T-1$}
    \State Draw
    $u_t\sim\mathcal{U}(0,1)$
\EndFor

\State \Return $\bm{u}$

\end{algorithmic}
\end{breakablealgorithm}

\subsection{$k$-body Pauli-product task}

This task probes multiplicative nonlinear processing across a temporal
depth $k$,

\begin{equation}
y_t=
\begin{cases}
\displaystyle\prod_{j=1}^{k}
\cos(\pi u_{t-j}), & t\ge k,\\[4pt]
0, & t<k.
\end{cases}
\label{eq:kbody}
\end{equation}

\begin{breakablealgorithm}
\caption{$k$-Body Pauli Product Benchmark}
\label{alg:kbody}
\begin{algorithmic}[1]

\Require Sequence length $T$, temporal order $k$,
random seed $s$

\Ensure Input sequence $\bm{u}$ and target sequence $\bm{y}$

\State
$\bm{u}\gets\Call{RandomInput}{T,s}$

\State
$\bm{y}\gets\bm{0}_{T}$

\For{$t=k,\ldots,T-1$}

    \State $p\gets1$

    \For{$j=1,\ldots,k$}

        \State
        $p\gets
        p\cos(\pi u_{t-j})$

    \EndFor

    \State
    $y_t\gets p$

\EndFor

\State \Return $\bm{u},\bm{y}$

\end{algorithmic}
\end{breakablealgorithm}

\subsection{Temporal parity task}

The continuous input is thresholded as
$b_t=\mathbb I(u_t>0.5)$. The target is the parity of the previous
$k$ bits,

\begin{equation}
p_t=b_{t-1}\oplus\cdots\oplus b_{t-k},
\qquad
y_t=
\begin{cases}
2p_t-1, & t\ge k,\\
0, & t<k.
\end{cases}
\label{eq:parity_target}
\end{equation}

\begin{breakablealgorithm}
\caption{Temporal Parity ($k$-XOR) Benchmark}
\label{alg:parity}
\begin{algorithmic}[1]

\Require Sequence length $T$, temporal order $k$,
random seed $s$

\Ensure Input sequence $\bm{u}$ and target sequence $\bm{y}$

\State
$\bm{u}\gets\Call{RandomInput}{T,s}$

\For{$t=0,\ldots,T-1$}

    \If{$u_t>0.5$}
        \State $b_t\gets1$
    \Else
        \State $b_t\gets0$
    \EndIf

\EndFor

\State
$\bm{y}\gets\bm{0}_{T}$

\For{$t=k,\ldots,T-1$}

    \State $x\gets0$

    \For{$j=1,\ldots,k$}

        \State
        $x\gets x\oplus b_{t-j}$

    \EndFor

    \State
    $y_t\gets2x-1$

\EndFor

\State \Return $\bm{u},\bm{y}$

\end{algorithmic}
\end{breakablealgorithm}

\subsection{Sparse Volterra task}

This benchmark probes sparse nonlinear interactions between delayed
samples,

\begin{equation}
y_t=
\begin{cases}
u_{t-1}u_{t-k},
& k\le5,\; t\ge\max(1,k),\\[3pt]
u_{t-1}u_{t-5}u_{t-k},
& k>5,\; t\ge k,\\[3pt]
0, & \text{otherwise}.
\end{cases}
\label{eq:sparse_volterra}
\end{equation}

\begin{breakablealgorithm}
\caption{Sparse Volterra Kernel Benchmark}
\label{alg:volterra}
\begin{algorithmic}[1]

\Require Sequence length $T$, temporal order $k$,
random seed $s$

\Ensure Input sequence $\bm{u}$ and target sequence $\bm{y}$

\State
$\bm{u}\gets\Call{RandomInput}{T,s}$

\State
$\bm{y}\gets\bm{0}_{T}$

\If{$k\leq5$}

    \For{$t=\max(1,k),\ldots,T-1$}

        \State
        $y_t\gets
        u_{t-1}u_{t-k}$

    \EndFor

\Else

    \For{$t=k,\ldots,T-1$}

        \State
        $y_t\gets
        u_{t-1}u_{t-5}u_{t-k}$

    \EndFor

\EndIf

\State \Return $\bm{u},\bm{y}$

\end{algorithmic}
\end{breakablealgorithm}

\subsection{Dynamical phase-correlation task}

The fourth benchmark uses ten alternating stochastic segments. Phase A
is a first-order autoregressive process,

\begin{equation}
u_\tau^{(A)}
=
0.3\,u_{\tau-1}^{(A)}
+
0.7\,\xi_\tau,
\qquad
\xi_\tau\sim\mathcal U(0,1),
\label{eq:phaseA}
\end{equation}

whereas Phase B uses a $k$-tap recurrence,

\begin{equation}
u_\tau^{(B)}
=
\sum_{j=0}^{k-1}
\frac{0.85}{k}u_{\tau-1-j}^{(B)}
+
0.15\,\xi_\tau ,
\qquad \tau\ge k,
\label{eq:phaseB}
\end{equation}

with early samples drawn uniformly and the resulting sequence clipped
to $[0,1]$. After centering
$\tilde u_t=u_t-\bar u$, the target is a normalized lag-$k$
correlation over a $20$-sample window,

\begin{equation}
y_t
=
\frac{
\displaystyle
\sum_{\tau=\ell_t}^{t}
\tilde u_\tau\tilde u_{\tau-k}
}{
(t-\ell_t+1)
[\Var(\bm u)+10^{-9}]
},
\qquad
\ell_t=\max(k,t-19),
\label{eq:phase_target}
\end{equation}

whenever $\ell_t\le t$, and zero otherwise. Thus, despite the
``phase'' terminology in the code, the target is a continuous
correlation statistic rather than a binary phase label.

\begin{breakablealgorithm}
\caption{Quantum Phase Recognition Benchmark}
\label{alg:quantum_phase}
\begin{algorithmic}[1]

\Require Sequence length $T$, correlation delay $k$,
random seed $s$

\Ensure Input sequence $\bm{u}$ and target sequence $\bm{y}$

\State Initialize pseudorandom generator using seed $s$

\State
$\alpha_A\gets0.3$

\State
$N_{\mathrm{seg}}\gets10$

\State
$L\gets\lfloor T/N_{\mathrm{seg}}\rfloor$

\State
$\bm{u}\gets\bm{0}_{T}$

\For{$r=0,\ldots,N_{\mathrm{seg}}-1$}

    \State
    $\ell\gets rL$

    \If{$r<N_{\mathrm{seg}}-1$}
        \State $h\gets\ell+L$
    \Else
        \State $h\gets T$
    \EndIf

    \State
    $M\gets h-\ell$

    \If{$r\bmod2=0$}

        \Comment{Phase A}

        \State
        $v_0\sim\mathcal{U}(0,1)$

        \For{$\tau=1,\ldots,M-1$}

            \State
            $\xi\sim\mathcal{U}(0,1)$

            \State
            $v_\tau
            \gets
            \alpha_Av_{\tau-1}
            +(1-\alpha_A)\xi$

        \EndFor

    \Else

        \Comment{Phase B}

        \For{$j=0,\ldots,k-1$}
            \State
            $c_j\gets0.85/k$
        \EndFor

        \For{$\tau=0,\ldots,M-1$}

            \State
            $\xi\sim\mathcal{U}(0,1)$

            \If{$\tau<k$}

                \State
                $v_\tau\gets\xi$

            \Else

                \State
                \[
                v_\tau
                \gets
                \sum_{j=0}^{k-1}
                c_jv_{\tau-1-j}
                +0.15\xi
                \]

            \EndIf

        \EndFor

        \For{$\tau=0,\ldots,M-1$}
            \State
            $v_\tau
            \gets
            \min(1,\max(0,v_\tau))$
        \EndFor

    \EndIf

    \For{$\tau=0,\ldots,M-1$}
        \State
        $u_{\ell+\tau}\gets v_\tau$
    \EndFor

\EndFor

\State Compute
$\bar{u}\gets
T^{-1}\sum_{t=0}^{T-1}u_t$

\For{$t=0,\ldots,T-1$}
    \State
    $\tilde{u}_t\gets u_t-\bar{u}$
\EndFor

\State
$\sigma_u^2
\gets
\operatorname{Var}(\bm{u})+10^{-9}$

\State
$W_Q\gets20$

\State
$\bm{y}\gets\bm{0}_{T}$

\For{$t=0,\ldots,T-1$}

    \State
    $\ell_t
    \gets
    \max(k,t-W_Q+1)$

    \If{$\ell_t\leq t$}

        \State $A\gets0$

        \For{$\tau=\ell_t,\ldots,t$}

            \State
            $A
            \gets
            A+
            \tilde{u}_{\tau}
            \tilde{u}_{\tau-k}$

        \EndFor

        \State
        \[
        y_t
        \gets
        \frac{A}
        {(t-\ell_t+1)\sigma_u^2}
        \]

    \EndIf

\EndFor

\State \Return $\bm{u},\bm{y}$

\end{algorithmic}
\end{breakablealgorithm}

Together, the four benchmarks span multiplicative nonlinearity,
nonlinearly separable temporal parity, sparse high-order memory, and
changes in temporal-correlation structure. The parameter $k$ changes
the effective temporal difficulty, but its precise role depends on the
benchmark.

% =====================================================================
\section{PCSR: Critical-Phase Extension and Shadow Readout}
\label{sec:pcsr}
% =====================================================================

PCSR keeps the hardware topology and temporal encoder of
Section~\ref{sec:backbone}. Its defining change is to replace every
full CZ interaction by a continuously tunable controlled-phase gate.
Therefore, PCSR should be interpreted as a one-parameter family built
on the baseline reservoir, not as a second unrelated architecture.

\subsection{Critical-phase reservoir evolution}

For every reservoir edge,

\begin{equation}
CP(2g)
=
\operatorname{diag}(1,1,1,e^{i2g}).
\label{eq:critical_cp}
\end{equation}

The two limiting cases are

\begin{equation}
g=0:\;CP(0)=I,
\qquad
g=\frac{\pi}{2}:\;CP(\pi)=CZ.
\label{eq:g_limits}
\end{equation}

Hence, $g=\pi/2$ recovers the interaction structure of the baseline
full-CZ reservoir, while intermediate $g$ values continuously tune the
interaction phase. PCSR experiments use fixed local biases
$\beta_i^{(z)}\sim\mathcal U(0,2\pi)$ and
$\beta_i^{(x)}\sim\mathcal U(0.3,0.7)$.

Let $U_{\mathrm{step}}(g)$ denote one topology-dependent controlled-phase
layer followed by these fixed local rotations. For $R$ repetitions,

\begin{equation}
U_{\mathcal R}(g)
=
[U_{\mathrm{step}}(g)]^R.
\label{eq:pcsr_repeated_unitary}
\end{equation}

Crucially, the temporal encoder is exactly the encoder already defined
in Eq.~\eqref{eq:shared_encoder}; it is not redefined for PCSR. The
state is

\begin{equation}
\ket{\psi_t(g)}
=
U_{\mathcal R}(g)
U_{\mathrm{enc}}(\bm w_t)
\ket{0}^{\otimes N}.
\label{eq:pcsr_state}
\end{equation}

\begin{breakablealgorithm}
\caption{Critical Reservoir Unitary Construction}
\label{alg:critical_unitary}
\small
\begin{algorithmic}[1]

\Require Reservoir topology
$\mathcal{T}=(\mathcal{Q},\mathcal{E}_{\mathcal{R}})$,
interaction parameter $g$,
rotation biases
$\bm{\beta}^{(z)}$ and $\bm{\beta}^{(x)}$,
fixed-layer flag $F$

\Ensure Reservoir-step unitary $U_{\mathrm{step}}$

\State $N\gets|\mathcal{Q}|$

\State Initialize an $N$-qubit quantum circuit

\State Partition $\mathcal{E}_{\mathcal{R}}$
into non-overlapping interaction moments

\For{each interaction moment $\mathcal{M}$}

    \For{each edge $(i,j)\in\mathcal{M}$}

        \State Apply
        $CP_{ij}(2g)$

    \EndFor

\EndFor

\If{$F=\mathrm{true}$}

    \For{$i=1,\ldots,N$}
        \State Apply
        $R_z(\beta_i^{(z)})$
    \EndFor

    \For{$i=1,\ldots,N$}
        \State Apply
        $R_x(\beta_i^{(x)})$
    \EndFor

\EndIf

\State Convert the circuit to the corresponding
big-endian unitary $U_{\mathrm{step}}$

\State \Return $U_{\mathrm{step}}$

\end{algorithmic}
\end{breakablealgorithm}

\begin{breakablealgorithm}
\caption{PCSR Reservoir Evolution}
\label{alg:pcsr_reservoir}
\small
\begin{algorithmic}[1]

\Require Topology $\mathcal{T}$,
input sequence $\bm{u}$,
interaction parameter $g$,
window size $W$,
reservoir repetitions $R$,
random seed $s$

\Ensure Reservoir states
$\{\ket{\psi_t}\}_{t=0}^{T-1}$

\State $N\gets|\mathcal{Q}|$,
$T\gets|\bm{u}|$

\State Initialize pseudorandom generator with seed $s$

\State Draw
$\beta_i^{(z)}\sim\mathcal{U}(0,2\pi)$

\State Draw
$\beta_i^{(x)}\sim\mathcal{U}(0.3,0.7)$

\State
$U_{\mathrm{step}}
\gets
\Call{CriticalReservoirUnitary}
{\mathcal{T},g,\bm{\beta}^{(z)},\bm{\beta}^{(x)}}$

\State
$U_{\mathcal{R}}
\gets
(U_{\mathrm{step}})^R$

\State Define
$\phi_w$ linearly over $[0.5,1]$

\State Define
$q(w)=w\bmod N$

\State
$\ket{\psi_0^{\mathrm{init}}}
\gets
\ket{0}^{\otimes N}$

\For{$t=0,\ldots,T-1$}

    \State
    $\ell\gets\max(0,t-W+1)$

    \State Construct a zero-padded window
    $\bm{w}_t$ containing
    $\{u_{\ell},\ldots,u_t\}$

    \State
    $\theta_i\gets0$,
    $i=1,\ldots,N$

    \For{$w=0,\ldots,W-1$}

        \State $i\gets q(w)$

        \State
        $\theta_i
        \gets
        \theta_i+\pi w_{t,w}\phi_w$

    \EndFor

    \State
    $U_{\mathrm{enc}}
    \gets
    \bigotimes_{i=1}^{N}R_y(\theta_i)$

    \State
    $\ket{\psi_t}
    \gets
    U_{\mathcal{R}}
    U_{\mathrm{enc}}
    \ket{\psi_0^{\mathrm{init}}}$

\EndFor

\State
\Return
$\{\ket{\psi_t}\}_{t=0}^{T-1}$

\end{algorithmic}
\end{breakablealgorithm}

\subsection{Dynamical diagnostics}

The phase scan is interpreted using state and operator diagnostics. For
a bipartition with
$N_A=\lfloor N/2\rfloor$ and $N_B=N-N_A$, reshape the state amplitudes
into
$\Psi\in\mathbb C^{2^{N_A}\times2^{N_B}}$. If
$\sigma_i$ are its singular values, the half-system entropy is

\begin{equation}
S_{\mathrm{half}}
=
-\sum_{i:p_i>10^{-12}}p_i\ln p_i,
\qquad
p_i=\sigma_i^2.
\label{eq:half_entropy}
\end{equation}

Operator entanglement is obtained by realigning the unitary across the
same bipartition, normalizing the squared singular values, and
computing the same Shannon/von-Neumann form,

\begin{equation}
S_{\mathrm{op}}
=
-\sum_{i:p_i>10^{-12}}p_i\ln p_i,
\qquad
p_i=
\frac{\sigma_i^2}{\sum_j\sigma_j^2}.
\label{eq:operator_entropy}
\end{equation}

The implementation also defines the single-step spectral form factor

\begin{equation}
K(U)
=
\frac{|\Tr(U)|^2}{d},
\qquad d=2^N.
\label{eq:sff}
\end{equation}

\begin{breakablealgorithm}
\caption{Half-System Entanglement Entropy}
\label{alg:half_entropy}
\small
\begin{algorithmic}[1]

\Require Quantum state $\ket{\psi}$,
number of qubits $N$

\Ensure Half-system entropy $S_{\mathrm{half}}$

\State
$N_A\gets\lfloor N/2\rfloor$

\State
$N_B\gets N-N_A$

\State Reshape $\ket{\psi}$ into
\[
\Psi
\in
\mathbb{C}^{2^{N_A}\times2^{N_B}}
\]

\State Compute singular values
$\{\sigma_i\}$ of $\Psi$

\For{each $\sigma_i$}

    \State
    $p_i\gets\sigma_i^2$

\EndFor

\State Remove all $p_i\leq10^{-12}$

\State
$S_{\mathrm{half}}
\gets
-\sum_i p_i\ln p_i$

\State
\Return $S_{\mathrm{half}}$

\end{algorithmic}
\end{breakablealgorithm}

\begin{breakablealgorithm}
\caption{Operator Entanglement}
\label{alg:operator_entanglement}
\small
\begin{algorithmic}[1]

\Require Reservoir unitary $U$,
even number of qubits $N$

\Ensure Operator entanglement $S_{\mathrm{op}}$

\State
$d_h\gets2^{N/2}$

\State Reshape
$U$ into a rank-$4$ tensor of dimensions
$d_h\times d_h\times d_h\times d_h$

\State Permute tensor indices according to
$(1,3,2,4)$

\State Reshape the result into
\[
\widetilde{U}
\in
\mathbb{C}^{d_h^2\times d_h^2}
\]

\State Compute singular values
$\{\sigma_i\}$ of $\widetilde{U}$

\State
$p_i
\gets
\sigma_i^2/
\sum_j\sigma_j^2$

\State Remove all $p_i\leq10^{-12}$

\State
$S_{\mathrm{op}}
\gets
-\sum_i p_i\ln p_i$

\State
\Return $S_{\mathrm{op}}$

\end{algorithmic}
\end{breakablealgorithm}

\begin{breakablealgorithm}
\caption{Spectral Form Factor}
\label{alg:sff}
\small
\begin{algorithmic}[1]

\Require Reservoir unitary $U$

\Ensure Spectral form factor $K$

\State
$d\gets$ dimension of $U$

\State
$\tau\gets\operatorname{Tr}(U)$

\State
$K\gets|\tau|^2/d$

\State
\Return $K$

\end{algorithmic}
\end{breakablealgorithm}

\subsection{From exact Pauli features to shadow features}

PCSR expands the observable representation from the baseline
$\{X_i,Y_i,Z_i,Z_iZ_j\}$ set to all same-axis one- and two-body
observables implemented by the code,

\begin{equation}
\mathcal P_2
=
\{X_i,Y_i,Z_i\}_{i=1}^{N}
\cup
\{X_iX_j,Y_iY_j,Z_iZ_j\}_{i<j}.
\label{eq:pcsr_pauli_set}
\end{equation}

The resulting feature dimension is

\begin{equation}
D_{\mathrm{PCSR}}
=
3N+3\binom{N}{2}.
\label{eq:pcsr_feature_dim}
\end{equation}

Algorithm~\ref{alg:exact_pauli} computes these expectations exactly from
a state vector.

\begin{breakablealgorithm}
\caption{Exact Pauli Expectation Extraction}
\label{alg:exact_pauli}
\small
\begin{algorithmic}[1]

\Require Quantum state $\ket{\psi}$,
number of qubits $N$,
maximum Pauli weight $w_{\max}$

\Ensure Observable labels $\mathcal{L}$ and
expectation vector $\bm{v}$

\State
$\mathcal{L}\gets[\;]$,
$\bm{v}\gets[\;]$

\For{$i=1,\ldots,N$}

    \State Compute $\langle Z_i\rangle$

    \State Append $Z_i$ and $\langle Z_i\rangle$

    \State Compute $\langle X_i\rangle$

    \State Append $X_i$ and $\langle X_i\rangle$

    \State Compute $\langle Y_i\rangle$

    \State Append $Y_i$ and $\langle Y_i\rangle$

\EndFor

\If{$w_{\max}\geq2$}

    \For{$i=1,\ldots,N$}

        \For{$j=i+1,\ldots,N$}

            \State Compute $\langle Z_iZ_j\rangle$

            \State Append $Z_iZ_j$ and its expectation

            \State Compute $\langle X_iX_j\rangle$

            \State Append $X_iX_j$ and its expectation

            \State Compute the implemented
            $Y_iY_j$ correlation estimator

            \State Append $Y_iY_j$ and its value

        \EndFor

    \EndFor

\EndIf

\State
\Return $\mathcal{L},\bm{v}$

\end{algorithmic}
\end{breakablealgorithm}

Two finite-measurement routes are then distinguished. The first is a
shot-noise surrogate used in ideal simulations: the exact expectation
of a Pauli string $P$ of weight $w(P)$ is perturbed as

\begin{equation}
\widetilde P
=
\langle P\rangle
+
\sqrt{\frac{3^{w(P)}}{M}}\,
\xi,
\qquad
\xi\sim\mathcal N(0,1).
\label{eq:shadow_noise}
\end{equation}

This routine emulates the shot dependence of the feature vector but
does not explicitly sample measurement bases.

\begin{breakablealgorithm}
\caption{Shadow-Feature Construction}
\label{alg:shadow_features}
\small
\begin{algorithmic}[1]

\Require Reservoir states
$\{\ket{\psi_t}\}_{t=0}^{T-1}$,
number of qubits $N$,
shot number $M$,
maximum Pauli weight $w_{\max}$,
random seed $s$

\Ensure Observable labels $\mathcal{L}$ and
feature matrix $\bm{X}_{\mathrm{sh}}$

\State Initialize pseudorandom generator with seed $s$

\State
$(\mathcal{L},\cdot)
\gets
\Call{ExactPauliExpectations}
{\ket{\psi_0},N,w_{\max}}$

\State Initialize
$\bm{X}_{\mathrm{sh}}$

\For{$t=0,\ldots,T-1$}

    \State
    $(\cdot,\bm{v}_t)
    \gets
    \Call{ExactPauliExpectations}
    {\ket{\psi_t},N,w_{\max}}$

    \If{$M>0$}

        \For{$k=1,\ldots,|\mathcal{L}|$}

            \State Determine Pauli weight
            $w_k$ of observable $\mathcal{L}_k$

            \State
            $\sigma_k
            \gets
            \sqrt{3^{w_k}/M}$

            \State Draw
            $\xi\sim\mathcal{N}(0,1)$

            \State
            $X_{\mathrm{sh}}[t,k]
            \gets
            v_{t,k}+\sigma_k\xi$

        \EndFor

    \Else

        \State
        $\bm{X}_{\mathrm{sh}}[t,:]
        \gets
        \bm{v}_t$

    \EndIf

\EndFor

\State
\Return
$\mathcal{L},\bm{X}_{\mathrm{sh}}$

\end{algorithmic}
\end{breakablealgorithm}

The second route explicitly samples random local bases. Each qubit
chooses one basis transformation from

\begin{equation}
B_i\in\{I,H,S^\dagger H\},
\end{equation}

a computational-basis outcome is sampled after
$U_B=\bigotimes_i B_i$, and matched Pauli estimators receive the
standard $3^{w(P)}$ inverse-measurement factor. This is the route used
to formulate the hardware-level shadow reconstruction.

\begin{breakablealgorithm}
\caption{Randomized Measurement Feature Estimation}
\label{alg:random_measurement}
\small
\begin{algorithmic}[1]

\Require State $\ket{\psi}$,
number of qubits $N$,
number of shots $M$,
random seed $s$

\Ensure Estimated Pauli feature vector
$\widehat{\bm{v}}$

\State Define local basis transformations
\[
B_0=I,\qquad
B_1=H,\qquad
B_2=S^\dagger H
\]

\State Construct all single-qubit observables
from basis labels $\{1,2,0\}$

\State Construct all two-qubit observable pairs
from the same basis labels

\State Initialize
$\widehat{\bm{v}}\gets\bm{0}$

\For{$m=1,\ldots,M$}

    \For{$i=1,\ldots,N$}

        \State Randomly choose
        $b_i\in\{0,1,2\}$

    \EndFor

    \State
    $U_B
    \gets
    \bigotimes_{i=1}^{N}B_{b_i}$

    \State
    $\bm{p}\gets
    |U_B\ket{\psi}|^2$

    \State Normalize $\bm{p}$

    \State Sample computational-basis outcome
    $z\sim\bm{p}$

    \For{each target Pauli observable $P_k$}

        \If{all required bases match $P_k$}

            \State
            $s_k\gets1$

            \For{each qubit $i$ in
            $\operatorname{supp}(P_k)$}

                \State
                $s_k\gets
                s_k(-1)^{z_i}$

            \EndFor

            \State
            $\widehat{v}_k
            \gets
            \widehat{v}_k+
            3^{w(P_k)}s_k$

        \EndIf

    \EndFor

\EndFor

\State
$\widehat{\bm{v}}
\gets
\widehat{\bm{v}}/M$

\State
\Return $\widehat{\bm{v}}$

\end{algorithmic}
\end{breakablealgorithm}

\subsection{LUQPI readout and classical comparison}

After a washout period $W_0$, the usable indices are shuffled and split
into training and test sets. This interleaved split is used by the
generic LUQPI evaluation routine.

\begin{breakablealgorithm}
\caption{Interleaved Train--Test Partition}
\label{alg:interleaved_split}
\small
\begin{algorithmic}[1]

\Require Sequence length $T$,
training fraction $\eta$,
washout $W_0$,
random seed $s$

\Ensure Training indices $\mathcal{I}_{\mathrm{tr}}$
and test indices $\mathcal{I}_{\mathrm{te}}$

\State
$\mathcal{I}
\gets
\{W_0,\ldots,T-1\}$

\State Initialize pseudorandom generator
with seed $s+7$

\State Randomly shuffle $\mathcal{I}$

\State
$N_{\mathrm{tr}}
\gets
\lfloor\eta|\mathcal{I}|\rfloor$

\State
$\mathcal{I}_{\mathrm{tr}}
\gets
\mathcal{I}[1:N_{\mathrm{tr}}]$

\State
$\mathcal{I}_{\mathrm{te}}
\gets
\mathcal{I}[N_{\mathrm{tr}}+1:|\mathcal{I}|]$

\State
\Return
$\mathcal{I}_{\mathrm{tr}},
\mathcal{I}_{\mathrm{te}}$

\end{algorithmic}
\end{breakablealgorithm}

The shadow-derived features are passed either to a two-layer MLP or to
Ridge regression. The error metric remains the NRMSE in
Eq.~\eqref{eq:nrmse}; it is not a new PCSR-specific metric.

\begin{breakablealgorithm}
\caption{LUQPI Training and Evaluation}
\label{alg:luqpi}
\small
\begin{algorithmic}[1]

\Require Shadow feature matrix
$\bm{X}_{\mathrm{sh}}$,
target $\bm{y}$,
model type $\mathcal{M}$,
regularization $\alpha$,
hidden architecture $\bm{h}$,
training fraction $\eta$,
washout $W_0$,
seed $s$

\Ensure NRMSE, trained model $f$, predictions
$\hat{\bm{y}}$

\State
$(\mathcal{I}_{\mathrm{tr}},
\mathcal{I}_{\mathrm{te}})
\gets
\Call{InterleavedSplit}
{T,\eta,W_0,s}$

\If{$\mathcal{M}=\mathrm{MLP}$}

    \State Initialize MLP with hidden layers $\bm{h}$,
    regularization $\alpha$, and maximum $500$ iterations

\Else

    \State Initialize Ridge regression
    with regularization $\alpha$

\EndIf

\State Fit
\[
f:
\bm{X}_{\mathrm{sh}}
[\mathcal{I}_{\mathrm{tr}}]
\rightarrow
\bm{y}
[\mathcal{I}_{\mathrm{tr}}]
\]

\State Predict
\[
\hat{\bm{y}}
\gets
f
\left(
\bm{X}_{\mathrm{sh}}
[\mathcal{I}_{\mathrm{te}}]
\right)
\]

\State
$\bm{y}_{\mathrm{te}}
\gets
\bm{y}[\mathcal{I}_{\mathrm{te}}]$

\State Compute
\[
\mathrm{NRMSE}
\gets
\sqrt{
\frac{
\operatorname{mean}
[
(\hat{\bm{y}}
-\bm{y}_{\mathrm{te}})^2
]
}{
\operatorname{Var}
(\bm{y}_{\mathrm{te}})
+10^{-12}
}
}
\]

\State
\Return
NRMSE, $f$, $\hat{\bm{y}}$

\end{algorithmic}
\end{breakablealgorithm}

To determine whether the quantum representation adds information
beyond direct access to the input history, a classical-only baseline is
constructed from $L$ lagged inputs,

\begin{equation}
\bm x_t^{(\mathrm C)}
=
[u_t,u_{t-1},\ldots,u_{t-L+1}],
\label{eq:classical_lags}
\end{equation}

with early entries zero padded.

\begin{breakablealgorithm}
\caption{Classical-Only Lagged Baseline}
\label{alg:classical_baseline}
\small
\begin{algorithmic}[1]

\Require Input sequence $\bm{u}$,
target sequence $\bm{y}$,
number of lags $L$,
model type $\mathcal{M}$,
hidden architecture $\bm{h}$,
training fraction $\eta$,
washout $W_0$,
regularization $\alpha$,
seed $s$

\Ensure Baseline NRMSE and trained model $f_{\mathrm{cl}}$

\State
$T\gets|\bm{u}|$

\State Initialize
\[
\bm{X}_{\mathrm{cl}}
\gets
\bm{0}_{T\times L}
\]

\For{$i=0,\ldots,L-1$}

    \For{$t=i,\ldots,T-1$}

        \State
        $X_{\mathrm{cl}}[t,i]
        \gets
        u_{t-i}$

    \EndFor

\EndFor

\State
$(\mathcal{I}_{\mathrm{tr}},
\mathcal{I}_{\mathrm{te}})
\gets
\Call{InterleavedSplit}
{T,\eta,W_0,s}$

\If{$\mathcal{M}=\mathrm{MLP}$}

    \State Initialize MLP with hidden layers $\bm{h}$,
    regularization $\alpha$, and maximum $500$ iterations

\Else

    \State Initialize Ridge regression
    with regularization $\alpha$

\EndIf

\State Fit
$f_{\mathrm{cl}}$ using
$\bm{X}_{\mathrm{cl}}
[\mathcal{I}_{\mathrm{tr}}]$
and
$\bm{y}[\mathcal{I}_{\mathrm{tr}}]$

\State
$\hat{\bm{y}}_{\mathrm{cl}}
\gets
f_{\mathrm{cl}}
(
\bm{X}_{\mathrm{cl}}
[\mathcal{I}_{\mathrm{te}}]
)$

\State Compute
\[
\mathrm{NRMSE}_{\mathrm{cl}}
=
\sqrt{
\frac{
\operatorname{mean}
[
(\hat{\bm{y}}_{\mathrm{cl}}
-\bm{y}_{\mathrm{te}})^2
]
}{
\operatorname{Var}
(\bm{y}_{\mathrm{te}})
+10^{-12}
}
}
\]

\State
\Return
$\mathrm{NRMSE}_{\mathrm{cl}},
f_{\mathrm{cl}}$

\end{algorithmic}
\end{breakablealgorithm}

The PCSR processing chain is therefore

\begin{equation}
\bm u
\rightarrow
U_{\mathrm{enc}}(\bm w_t)
\rightarrow
U_{\mathcal R}(g)
\rightarrow
\ket{\psi_t(g)}
\rightarrow
\bm X_{\mathrm Q}
\rightarrow
f_{\mathrm{LUQPI}}
\rightarrow
\hat{\bm y},
\label{eq:pcsr_pipeline}
\end{equation}

where $\bm X_{\mathrm Q}$ denotes the exact or shadow-derived quantum
feature matrix. The experiments below vary only selected parts of this
chain while keeping the common definitions fixed.

% =====================================================================
\section{Experimental Protocol}
\label{sec:experiments}
% =====================================================================

The experiments are designed to answer three progressively stronger
questions. First, does changing $g$ alter both quantum-information
diagnostics and task performance? Second, do PCSR features add
predictive information beyond a classical lag vector, and is that
advantage stable across training sizes and random seeds? Third, does
the representation remain useful under finite-shot noise and
hardware-realistic execution?

\subsection{Interaction-phase scan}

For the PCSR-specific memory analysis, delayed targets are formed on
the same interleaved train/test indices used by the LUQPI routines and
Eq.~\eqref{eq:mc} is evaluated for each delay. The total memory
reported in the scan is

\begin{equation}
MC_{\mathrm{tot}}
=
\sum_{k=1}^{K_{\max}}MC_k.
\label{eq:mc_total}
\end{equation}

\begin{breakablealgorithm}
\caption{Delay-Dependent PCSR Memory Capacity}
\label{alg:pcsr_memory}
\small
\begin{algorithmic}[1]

\Require PCSR feature matrix $\bm{X}$,
input sequence $\bm{u}$,
maximum delay $K_{\max}$,
training fraction $\eta$,
washout $W_0$

\Ensure Memory coefficients
$\{MC_k\}_{k=1}^{K_{\max}}$

\State $T\gets|\bm{u}|$

\State
$(\mathcal{I}_{\mathrm{tr}},
\mathcal{I}_{\mathrm{te}})
\gets
\Call{InterleavedSplit}
{T,\eta,W_0}$

\State Initialize
$MC_k\gets0$,
$k=1,\ldots,K_{\max}$

\For{$k=1,\ldots,K_{\max}$}

    \If{$k\geq T-W_0-5$}
        \State \textbf{break}
    \EndIf

    \State Initialize
    $\bm{d}^{(k)}\gets\bm{0}_T$

    \For{$t=k,\ldots,T-1$}
        \State
        $d_t^{(k)}\gets u_{t-k}$
    \EndFor

    \State Initialize Ridge regression with
    $\alpha=10^{-6}$

    \State Fit the model using
    \[
    \bm{X}[\mathcal{I}_{\mathrm{tr}}]
    \rightarrow
    \bm{d}^{(k)}[\mathcal{I}_{\mathrm{tr}}]
    \]

    \State Predict
    \[
    \hat{\bm{d}}^{(k)}
    \gets
    f_k
    \left(
    \bm{X}[\mathcal{I}_{\mathrm{te}}]
    \right)
    \]

    \State
    $\bm{d}^{(k)}_{\mathrm{te}}
    \gets
    \bm{d}^{(k)}
    [\mathcal{I}_{\mathrm{te}}]$

    \State Compute
    \[
    MC_k
    \gets
    \frac{
    \Cov
    \left(
    \hat{\bm{d}}^{(k)},
    \bm{d}^{(k)}_{\mathrm{te}}
    \right)^2
    }{
    \Var
    \left(
    \hat{\bm{d}}^{(k)}
    \right)
    \Var
    \left(
    \bm{d}^{(k)}_{\mathrm{te}}
    \right)
    +10^{-12}
    }
    \]

\EndFor

\State
\Return
$\{MC_k\}_{k=1}^{K_{\max}}$

\end{algorithmic}
\end{breakablealgorithm}

The interaction strength is scanned over
$g\in[0,\pi/2]$, using eleven equally spaced points when no custom grid
is supplied. The same topology and local biases are retained across
the scan. For each $g$, the experiment records mean half-system entropy
over the final 30 random-input states, operator entanglement of
$U_{\mathrm{step}}^2$, total memory capacity, and task NRMSE for the
$k$-body Pauli and parity benchmarks. This directly links the
interaction regime to both physical and computational observables.

\begin{breakablealgorithm}
\caption{PCSR Interaction-Phase Scan}
\label{alg:phase_scan}
\small
\begin{algorithmic}[1]

\Require Topology type $\mathcal{T}$,
number of qubits $N$,
sequence length $T$,
interaction values $\mathcal{G}$,
$k$-body order $k_{\mathrm{P}}$,
seed $s$

\Ensure Entanglement, memory, and task-performance
metrics as functions of $g$

\If{$\mathcal{G}$ is not specified}
    \State
    $\mathcal{G}\gets
    \operatorname{linspace}(0,\pi/2,11)$
\EndIf

\State
$\mathcal{R}\gets
\Call{BuildTopology}{\mathcal{T},N}$

\State Initialize pseudorandom generator with seed $s$

\State Draw fixed
$\beta_i^{(z)}\sim\mathcal{U}(0,2\pi)$

\State Draw fixed
$\beta_i^{(x)}\sim\mathcal{U}(0.3,0.7)$

\State Generate $k$-body benchmark
\[
(\bm{u}_{\mathrm{kP}},\bm{y}_{\mathrm{kP}})
\gets
\Call{KBodyPauli}{T,k_{\mathrm{P}},s}
\]

\State Generate parity benchmark with $k=3$
\[
(\bm{u}_{\mathrm{par}},\bm{y}_{\mathrm{par}})
\gets
\Call{TemporalParity}{T,3,s}
\]

\State Generate independent random sequence
\[
\bm{u}_{\mathrm{rand}}
\sim
\mathcal{U}(0,1)^T
\]
using seed $s+1$

\For{each $g\in\mathcal{G}$}

    \State
    $\Psi_{\mathrm{kP}}
    \gets
    \Call{PCSRReservoir}
    {\mathcal{R},\bm{u}_{\mathrm{kP}},g}$

    \State
    $\Psi_{\mathrm{par}}
    \gets
    \Call{PCSRReservoir}
    {\mathcal{R},\bm{u}_{\mathrm{par}},g}$

    \State
    $\Psi_{\mathrm{rand}}
    \gets
    \Call{PCSRReservoir}
    {\mathcal{R},\bm{u}_{\mathrm{rand}},g}$

    \State Compute
    \[
    \overline{S}_{1/2}
    \gets
    \frac{1}{30}
    \sum_{t=T-30}^{T-1}
    S_{1/2}
    (\Psi_{\mathrm{rand},t})
    \]

    \State
    $U
    \gets
    \Call{CriticalReservoirUnitary}
    {\mathcal{R},g,
    \bm{\beta}^{(z)},\bm{\beta}^{(x)}}$

    \State
    $S_{\mathrm{op}}
    \gets
    \Call{OperatorEntanglement}{U^2,N}$

    \State
    $\bm{X}_{\mathrm{kP}}
    \gets
    \Call{ShadowFeatures}
    {\Psi_{\mathrm{kP}},N,M=0}$

    \State
    $\bm{X}_{\mathrm{par}}
    \gets
    \Call{ShadowFeatures}
    {\Psi_{\mathrm{par}},N,M=0}$

    \State
    $\bm{X}_{\mathrm{rand}}
    \gets
    \Call{ShadowFeatures}
    {\Psi_{\mathrm{rand}},N,M=0}$

    \State Train Ridge readout on
    $(\bm{X}_{\mathrm{kP}},
    \bm{y}_{\mathrm{kP}})$
    with $\alpha=10^{-4}$

    \State Record
    $\mathrm{NRMSE}_{\mathrm{kP}}$

    \State Train MLP readout with hidden layers
    $(64,32)$ on
    $(\bm{X}_{\mathrm{par}},
    \bm{y}_{\mathrm{par}})$

    \State Record
    $\mathrm{NRMSE}_{\mathrm{par}}$

    \State
    $\{MC_k\}
    \gets
    \Call{PCSRCMemory}
    {\bm{X}_{\mathrm{rand}},
    \bm{u}_{\mathrm{rand}},10}$

    \State
    $MC_{\mathrm{tot}}
    \gets
    \sum_{k=1}^{10}MC_k$

    \State Store
    \[
    \left(
    g,
    \overline{S}_{1/2},
    S_{\mathrm{op}},
    MC_{\mathrm{tot}},
    \mathrm{NRMSE}_{\mathrm{kP}},
    \mathrm{NRMSE}_{\mathrm{par}}
    \right)
    \]

\EndFor

\State \Return all phase-scan results

\end{algorithmic}
\end{breakablealgorithm}

\subsection{Classical--quantum representation separation}

Three representations are compared on exactly the same training and
test indices:

\begin{equation}
\bm X_{\mathrm C}
=
\text{eight-lag classical input features},
\qquad
\bm X_{\mathrm Q}
=
\text{PCSR shadow features},
\end{equation}

and

\begin{equation}
\bm X_{\mathrm{CQ}}
=
[\bm X_{\mathrm Q},\bm X_{\mathrm C}].
\label{eq:combined_features}
\end{equation}

The final 100 samples are held out for testing, while training examples
are sampled from the post-washout region preceding that test block.
The single-seed sweep measures sample efficiency as the training size
changes.

\begin{breakablealgorithm}
\caption{LUQPI Representation-Separation Sweep}
\label{alg:luqpi_separation}
\small
\begin{algorithmic}[1]

\Require Topology $\mathcal{T}$,
qubit number $N$,
sequence length $T$,
interaction strength $g$,
training sizes $\mathcal{N}_{\mathrm{tr}}$,
task set $\mathcal{B}$,
seed $s$

\Ensure NRMSE versus training size for classical,
shadow, and combined representations

\State
$\mathcal{R}
\gets
\Call{BuildTopology}{\mathcal{T},N}$

\State Draw fixed biases
$\bm{\beta}^{(z)}$ and
$\bm{\beta}^{(x)}$

\For{each benchmark $b\in\mathcal{B}$}

    \State Select benchmark parameters:
    \[
    \begin{array}{ll}
    \mathrm{kPauli}: & k=3,\\
    \mathrm{Parity}: & k=3,\\
    \mathrm{Volterra}: & k=4,\\
    \mathrm{QPhase}: & k=4
    \end{array}
    \]

    \State Generate
    $(\bm{u},\bm{y})$
    for benchmark $b$

    \State
    $\Psi
    \gets
    \Call{PCSRReservoir}
    {\mathcal{R},\bm{u},g}$

    \State
    $\bm{X}_{\mathrm{Q}}
    \gets
    \Call{ShadowFeatures}{\Psi,N,M=0}$

    \State Initialize
    $\bm{X}_{\mathrm{C}}
    \gets
    \bm{0}_{T\times8}$

    \For{$i=0,\ldots,7$}
        \For{$t=i,\ldots,T-1$}
            \State
            $X_{\mathrm{C}}[t,i]
            \gets u_{t-i}$
        \EndFor
    \EndFor

    \State
    $\bm{X}_{\mathrm{CQ}}
    \gets
    [
    \bm{X}_{\mathrm{Q}},
    \bm{X}_{\mathrm{C}}
    ]$

    \For{each training size
    $n_{\mathrm{tr}}\in
    \mathcal{N}_{\mathrm{tr}}$}

        \State
        $\mathcal{I}_{\mathrm{pool}}
        \gets
        \{30,\ldots,T-101\}$

        \State Shuffle
        $\mathcal{I}_{\mathrm{pool}}$
        using seed $s+7$

        \State
        $\mathcal{I}_{\mathrm{tr}}
        \gets$
        first $n_{\mathrm{tr}}$ shuffled indices

        \State
        $\mathcal{I}_{\mathrm{te}}
        \gets
        \{T-100,\ldots,T-1\}$

        \For{each
        $\bm{X}\in
        \{
        \bm{X}_{\mathrm{C}},
        \bm{X}_{\mathrm{Q}},
        \bm{X}_{\mathrm{CQ}}
        \}$}

            \State Initialize MLP with hidden layers
            $(64,32)$,
            $\alpha=10^{-3}$,
            and $400$ maximum iterations

            \State Fit using
            $\bm{X}[\mathcal{I}_{\mathrm{tr}}]$
            and
            $\bm{y}[\mathcal{I}_{\mathrm{tr}}]$

            \State Predict on
            $\mathcal{I}_{\mathrm{te}}$

            \State Compute and store
            \[
            \mathrm{NRMSE}
            =
            \sqrt{
            \frac{
            \operatorname{mean}
            [(\hat{\bm{y}}-\bm{y}_{\mathrm{te}})^2]
            }{
            \Var(\bm{y}_{\mathrm{te}})
            +10^{-12}
            }
            }
            \]

        \EndFor

    \EndFor

\EndFor

\State \Return all NRMSE curves

\end{algorithmic}
\end{breakablealgorithm}

The same separation test is then repeated over multiple independent
seeds. For a fixed seed and training size, all three representations
use the same shuffled training indices, making the comparison paired.
If $e_{m,s}(n)$ is the NRMSE of method $m$ at training size $n$ and
seed $s$, the reported statistics are

\begin{equation}
\bar e_m(n)
=
\frac{1}{S}\sum_{s=1}^{S} e_{m,s}(n),
\qquad
\sigma_m(n)
=
\sqrt{
\frac{1}{S}
\sum_{s=1}^{S}
[e_{m,s}(n)-\bar e_m(n)]^2
}.
\label{eq:multiseed_stats}
\end{equation}

\begin{breakablealgorithm}
\caption{Multi-Seed LUQPI Separation Analysis}
\label{alg:multiseed_luqpi}
\small
\begin{algorithmic}[1]

\Require Benchmark task $\mathcal{B}$,
number of qubits $N$,
sequence length $T$,
interaction strength $g$,
training sizes $\mathcal{N}_{\mathrm{tr}}$,
number of seeds $S$,
task order $k$,
topology $\mathcal{T}$

\Ensure Mean and standard deviation of NRMSE for
classical, shadow, and combined representations

\State Initialize result sets
\[
\mathcal{R}_{\mathrm{C}},
\mathcal{R}_{\mathrm{Q}},
\mathcal{R}_{\mathrm{CQ}}
\gets [\;]
\]

\For{$s=0,\ldots,S-1$}

    \State
    $\mathcal{R}
    \gets
    \Call{BuildTopology}{\mathcal{T},N}$

    \State Initialize reservoir random generator
    using seed $s+100$

    \For{$i=1,\ldots,N$}
        \State
        $\beta_i^{(z)}
        \sim
        \mathcal{U}(0,2\pi)$
        \State
        $\beta_i^{(x)}
        \sim
        \mathcal{U}(0.3,0.7)$
    \EndFor

    \State Generate benchmark sequence
    \[
    (\mathbf{u},\mathbf{y})
    \gets
    \Call{BenchmarkTask}{\mathcal{B},T,k,s}
    \]

    \State Generate PCSR reservoir states
    \[
    \Psi
    \gets
    \Call{PCSRReservoir}
    {
    \mathcal{R},
    \mathbf{u},
    g,
    \boldsymbol{\beta}^{(z)},
    \boldsymbol{\beta}^{(x)}
    }
    \]

    \State Extract exact shadow representation
    \[
    \mathbf{X}_{\mathrm{Q}}
    \gets
    \Call{ShadowFeatures}{\Psi,N,M=0}
    \]

    \State Initialize
    \[
    \mathbf{X}_{\mathrm{C}}
    \gets
    \mathbf{0}_{T\times 8}
    \]

    \For{$i=0,\ldots,7$}

        \For{$t=i,\ldots,T-1$}

            \State
            $X_{\mathrm{C}}[t,i]
            \gets
            u_{t-i}$

        \EndFor

    \EndFor

    \State Construct combined representation
    \[
    \mathbf{X}_{\mathrm{CQ}}
    \gets
    [
    \mathbf{X}_{\mathrm{Q}},
    \mathbf{X}_{\mathrm{C}}
    ]
    \]

    \State Initialize per-seed NRMSE arrays
    $\mathbf{r}_{\mathrm{C}}$,
    $\mathbf{r}_{\mathrm{Q}}$,
    and
    $\mathbf{r}_{\mathrm{CQ}}$

    \For{each training size
    $n_{\mathrm{tr}}\in\mathcal{N}_{\mathrm{tr}}$}

        \State Set washout
        $W_0\gets30$

        \State Define fixed test indices
        \[
        \mathcal{I}_{\mathrm{te}}
        \gets
        \{T-100,\ldots,T-1\}
        \]

        \State Define training pool
        \[
        \mathcal{I}_{\mathrm{pool}}
        \gets
        \{W_0,\ldots,T-101\}
        \]

        \State Shuffle
        $\mathcal{I}_{\mathrm{pool}}$
        using seed $s+7$

        \State
        $\mathcal{I}_{\mathrm{tr}}
        \gets$
        first $n_{\mathrm{tr}}$ shuffled indices

        \For{each representation
        $\mathbf{X}\in
        \{
        \mathbf{X}_{\mathrm{C}},
        \mathbf{X}_{\mathrm{Q}},
        \mathbf{X}_{\mathrm{CQ}}
        \}$}

            \State Initialize MLP with hidden layers
            $(64,32)$,
            regularization
            $\alpha=10^{-3}$,
            and maximum $500$ iterations

            \State Fit
            \[
            f:
            \mathbf{X}[\mathcal{I}_{\mathrm{tr}}]
            \rightarrow
            \mathbf{y}[\mathcal{I}_{\mathrm{tr}}]
            \]

            \State Predict
            \[
            \hat{\mathbf{y}}
            \gets
            f
            \left(
            \mathbf{X}[\mathcal{I}_{\mathrm{te}}]
            \right)
            \]

            \State Compute
            \[
            \mathrm{NRMSE}
            \gets
            \sqrt{
            \frac{
            \operatorname{mean}
            \left[
            (\hat{\mathbf{y}}
            -
            \mathbf{y}_{\mathrm{te}})^2
            \right]
            }{
            \operatorname{Var}
            (\mathbf{y}_{\mathrm{te}})
            +10^{-12}
            }
            }
            \]

            \State Append NRMSE to the corresponding
            representation array

        \EndFor

    \EndFor

    \State Append
    $\mathbf{r}_{\mathrm{C}}$,
    $\mathbf{r}_{\mathrm{Q}}$,
    and
    $\mathbf{r}_{\mathrm{CQ}}$
    to their multi-seed result matrices

\EndFor

\For{each method
$m\in
\{\mathrm{C},\mathrm{Q},\mathrm{CQ}\}$}

    \State Compute training-size-wise mean
    \[
    \mu_m(n)
    =
    \frac{1}{S}
    \sum_{s=1}^{S}
    R_{m,s}(n)
    \]

    \State Compute training-size-wise standard deviation
    \[
    \sigma_m(n)
    =
    \sqrt{
    \frac{1}{S}
    \sum_{s=1}^{S}
    \left(
    R_{m,s}(n)-\mu_m(n)
    \right)^2
    }
    \]

\EndFor

\State \Return
$\{
\mu_{\mathrm{C}},
\sigma_{\mathrm{C}},
\mu_{\mathrm{Q}},
\sigma_{\mathrm{Q}},
\mu_{\mathrm{CQ}},
\sigma_{\mathrm{CQ}}
\}$

\end{algorithmic}
\end{breakablealgorithm}

\subsection{Finite-shot robustness}

The finite-shot experiment reuses the same PCSR states but changes the
feature-extraction noise level. Exact expectations ($M=0$) define the
reference, and shot-noise-surrogate features are evaluated at
$M\in\{50,100,250,500,1000,5000\}$. The tested tasks are
$k$-body Pauli ($k=3$), parity ($k=3$), and sparse Volterra ($k=4$).

\begin{breakablealgorithm}
\caption{Shadow Shot-Robustness Evaluation}
\label{alg:shot_robustness}
\small
\begin{algorithmic}[1]

\Require Topology $\mathcal{T}$,
qubit number $N$,
sequence length $T$,
interaction parameter $g$,
shot levels $\mathcal{K}$,
seed $s$

\Ensure Exact and finite-shot NRMSE for each task

\State
$\mathcal{R}
\gets
\Call{BuildTopology}{\mathcal{T},N}$

\State Draw fixed biases
$\bm{\beta}^{(z)}$ and
$\bm{\beta}^{(x)}$

\State Define benchmark set
\[
\mathcal{B}
=
\{
(\mathrm{kPauli},3),
(\mathrm{Parity},3),
(\mathrm{Volterra},4)
\}
\]

\For{each $(b,k)\in\mathcal{B}$}

    \State Generate benchmark
    $(\bm{u},\bm{y})$

    \State
    $\Psi
    \gets
    \Call{PCSRReservoir}
    {\mathcal{R},\bm{u},g}$

    \State Obtain exact features
    \[
    \bm{X}_{\mathrm{exact}}
    \gets
    \Call{ShadowFeatures}
    {\Psi,N,M=0}
    \]

    \State Train MLP with hidden layers $(64,32)$
    using $\bm{X}_{\mathrm{exact}}$

    \State Store exact-reference NRMSE
    $\mathrm{NRMSE}_{\mathrm{exact}}$

    \For{each shot level $K\in\mathcal{K}$}

        \State Generate noisy shadow features
        \[
        \bm{X}_K
        \gets
        \Call{ShadowFeatures}
        {\Psi,N,M=K,s+K}
        \]

        \State Train MLP with hidden layers $(64,32)$
        using $\bm{X}_K$

        \State Compute
        $\mathrm{NRMSE}_K$

        \State Store
        $\mathrm{NRMSE}_K$

    \EndFor

\EndFor

\State
\Return exact and finite-shot NRMSE results

\end{algorithmic}
\end{breakablealgorithm}

\subsection{Inference latency}

Latency is reported for four deployment situations:
(i) classical MLP prediction on precomputed lag features,
(ii) PCSR MLP prediction on precomputed shadow features,
(iii) an end-to-end state-vector reservoir path that includes window
encoding, state evolution, Pauli extraction, and MLP prediction, and
(iv) a real-QPU estimate based on an assumed circuit-execution rate.

The first two quantities are therefore \emph{cached-feature readout
latencies}, whereas the third includes reservoir computation and is not
directly the same timing quantity. The real-QPU estimate is

\begin{equation}
\tau_{\mathrm{QPU}}
=
\frac{M_{\mathrm{shot}}}
{R_{\mathrm{QPU}}},
\qquad
R_{\mathrm{QPU}}=5000~\mathrm{s}^{-1},
\label{eq:qpu_latency}
\end{equation}

where $R_{\mathrm{QPU}}$ is an assumed implementation constant rather
than a measured device rate in this routine.

\begin{breakablealgorithm}
\caption{Inference-Latency Benchmark}
\label{alg:latency}
\small
\begin{algorithmic}[1]

\Require Number of qubits $N$,
number of timed predictions $M_{\mathrm{test}}$,
QPU shots $M_{\mathrm{shot}}$,
seed $s$

\Ensure Latencies for classical-only, cached-PCSR,
simulated reservoir, and estimated real-QPU inference

\State Construct Chain topology with $N$ qubits

\State Draw fixed reservoir biases

\State Set
$g\gets\pi/4$

\State Generate $k$-body training task with
$T=300$ and $k=3$

\State Generate PCSR states for the training input

\State Extract exact shadow features
$\bm{X}_{\mathrm{Q,tr}}$

\State Train PCSR MLP with hidden layers $(64,32)$
after a washout of $30$

\State Construct eight-lag classical features
$\bm{X}_{\mathrm{C,tr}}$

\State Train classical MLP with hidden layers $(64,32)$
after a washout of $30$

\State Generate random test sequence
$\bm{u}_{\mathrm{te}}$

\State Precompute PCSR states and exact shadow features
$\bm{X}_{\mathrm{Q,te}}$

\State Precompute classical lag features
$\bm{X}_{\mathrm{C,te}}$

\State Start timer

\For{$m=1,\ldots,M_{\mathrm{test}}$}
    \State Evaluate classical MLP on one
    precomputed classical feature vector
\EndFor

\State Compute
\[
\tau_{\mathrm{C}}
\gets
\frac{\Delta t_{\mathrm{C}}}
{M_{\mathrm{test}}}\times1000
\]

\State Start timer

\For{$m=1,\ldots,M_{\mathrm{test}}$}
    \State Evaluate PCSR MLP on one
    precomputed shadow feature vector
\EndFor

\State Compute
\[
\tau_{\mathrm{PCSR}}
\gets
\frac{\Delta t_{\mathrm{PCSR}}}
{M_{\mathrm{test}}}\times1000
\]

\State Construct
\[
U_{\mathrm{step}}
\gets
\left[
\Call{CriticalReservoirUnitary}
{\mathcal{R},g}
\right]^2
\]

\State Set window size $W=8$

\State
$M_{\mathrm{sim}}
\gets
\min(M_{\mathrm{test}},20)$

\State Start timer

\For{$m=0,\ldots,M_{\mathrm{sim}}-1$}

    \State Construct the current zero-padded
    temporal window

    \State Construct
    $U_{\mathrm{enc}}$

    \State Simulate
    \[
    \ket{\psi_m}
    \gets
    U_{\mathrm{step}}
    U_{\mathrm{enc}}
    \ket{0}^{\otimes N}
    \]

    \State Extract exact one- and two-body
    Pauli features $\bm{x}_m$

    \State Evaluate PCSR MLP using $\bm{x}_m$

\EndFor

\State Compute
\[
\tau_{\mathrm{sim}}
\gets
\frac{\Delta t_{\mathrm{sim}}}
{M_{\mathrm{sim}}}\times1000
\]

\State Set assumed real-QPU execution rate
\[
R_{\mathrm{QPU}}\gets5000~\mathrm{s}^{-1}
\]

\State Estimate
\[
\tau_{\mathrm{QPU}}
\gets
\frac{M_{\mathrm{shot}}}
{R_{\mathrm{QPU}}}
\times1000
\]

\State \Return
$\tau_{\mathrm{C}}$,
$\tau_{\mathrm{PCSR}}$,
$\tau_{\mathrm{sim}}$,
$\tau_{\mathrm{QPU}}$

\end{algorithmic}
\end{breakablealgorithm}
% =====================================================================
\section{IBM Heron Noisy-Hardware Validation}
\label{sec:ibm_validation}
% =====================================================================

The final stage validates the same PCSR construction under a calibrated
IBM Heron execution model. When real-hardware execution is disabled,
the sampler is backed by an Aer simulator instantiated from the backend
snapshot; when enabled, the runtime sampler is bound to the physical
backend. This section therefore changes the execution and measurement
model, not the mathematical definition of the PCSR reservoir.

\subsection{Backend and hardware-level circuit}

\begin{breakablealgorithm}
\caption{IBM Heron Noisy Sampler Initialization}
\label{alg:ibm_sampler}
\small
\begin{algorithmic}[1]

\Require IBM backend $B$ and
hardware-execution flag $F_{\mathrm{real}}$

\Ensure Configured sampler $\mathcal{S}$

\If{a sampler is already cached}
    \State \Return cached sampler
\EndIf

\State
$B\gets\Call{GetIBMBackend}{}$

\If{$F_{\mathrm{real}}=\mathrm{true}$}

    \State
    $\mathcal{S}
    \gets
    \Call{SamplerV2}{B}$

    \State Store execution backend as $B$

\Else

    \State Construct hardware-calibrated noisy simulator
    \[
    B_{\mathrm{sim}}
    \gets
    \Call{AerSimulator.from\_backend}{B}
    \]

    \State
    $\mathcal{S}
    \gets
    \Call{SamplerV2}{B_{\mathrm{sim}}}$

    \State Store original backend $B$ as the
    compilation target

\EndIf

\State Cache $\mathcal{S}$

\State \Return $\mathcal{S}$

\end{algorithmic}
\end{breakablealgorithm}

For each time step, the hardware circuit implements the same
zero-padded window and the same $R_y$ encoder from
Eqs.~\eqref{eq:shared_window}--\eqref{eq:shared_encoder}. It then applies
the topology-dependent $CP(2g)$ interactions and fixed
$R_z$/$R_x$ rotations directly on the selected physical qubits. A
random local measurement basis is appended, with $I$, $H$, and
$S^\dagger H$ used for the implemented $Z$, $X$, and $Y$ settings,
respectively. The complete circuit is optionally transpiled against the
Heron backend.

\begin{breakablealgorithm}
\caption{Hardware-Level PCSR Circuit Construction}
\label{alg:build_ibm_pcsr}
\small
\begin{algorithmic}[1]

\Require Reservoir topology $\mathcal{T}$,
input window $\mathbf{w}$,
interaction strength $g$,
fixed biases
$\boldsymbol{\beta}^{(z)}$ and
$\boldsymbol{\beta}^{(x)}$,
measurement basis
$\mathbf{b}$,
measurement flag $F_{\mathrm{m}}$,
native-compilation flag $F_{\mathrm{c}}$

\Ensure Hardware-compatible PCSR circuit $C$

\State Obtain IBM backend $B$

\State
$N\gets|\mathcal{Q}|$

\State
$W\gets|\mathbf{w}|$

\State Define linearly spaced weights
\[
\phi_r\in[0.5,1],
\qquad
r=0,\ldots,W-1
\]

\State Define
\[
q(r)=r\bmod N
\]

\State Initialize
$\theta_i\gets0$,
$i=1,\ldots,N$

\For{$r=0,\ldots,W-1$}

    \State
    $i\gets q(r)$

    \State
    $\theta_i
    \gets
    \theta_i+\pi w_r\phi_r$

\EndFor

\State Initialize a circuit on all physical device qubits
and $N$ classical bits

\For{$i=1,\ldots,N$}

    \State Apply
    $R_y(\theta_i)$
    to physical reservoir qubit $q_i$

\EndFor

\State Partition topology edges into non-overlapping
interaction moments

\For{each interaction edge $(a,b)$}

    \State Apply
    \[
    CP_{a,b}(2g)
    \]

\EndFor

\For{$i=1,\ldots,N$}

    \State Apply
    $R_z(\beta_i^{(z)})$
    to physical qubit $q_i$

\EndFor

\For{$i=1,\ldots,N$}

    \State Apply
    $R_x(\beta_i^{(x)})$
    to physical qubit $q_i$

\EndFor

\For{$i=1,\ldots,N$}

    \If{$b_i=X$}

        \State Apply $H$ to $q_i$

    \ElsIf{$b_i=Y$}

        \State Apply $S^\dagger$ followed by $H$
        to $q_i$

    \Else

        \State Apply no basis rotation

    \EndIf

\EndFor

\If{$F_{\mathrm{m}}=\mathrm{true}$}

    \For{$i=1,\ldots,N$}

        \State Measure physical reservoir qubit
        $q_i$ into classical bit $i$

    \EndFor

\EndIf

\If{$F_{\mathrm{c}}=\mathrm{true}$}

    \State Transpile $C$ against backend $B$
    using optimization level $1$

\EndIf

\State \Return $C$

\end{algorithmic}
\end{breakablealgorithm}

\subsection{Measurement decoding and shadow reconstruction}

Sampler bit strings are decoded into an
$M\times N$ binary matrix, with column $i$ corresponding to reservoir
qubit $i$.

\begin{breakablealgorithm}
\caption{IBM Measurement Bitstring Decoding}
\label{alg:ibm_bit_decode}
\small
\begin{algorithmic}[1]

\Require Sampler result $R$ and
number of reservoir qubits $N$

\Ensure Measurement matrix
$\mathbf{Z}\in\{0,1\}^{M\times N}$

\State Extract classical-register measurement data
from $R$

\State Obtain all returned bit strings
$\{s_m\}_{m=1}^{M}$

\For{$m=1,\ldots,M$}

    \State Reverse bit string $s_m$

    \For{$i=1,\ldots,N$}

        \State
        $Z_{m,i}
        \gets$
        integer value of the $i$th reversed bit

    \EndFor

\EndFor

\State \Return $\mathbf{Z}$

\end{algorithmic}
\end{breakablealgorithm}

Across randomized basis circuits, the total number of shots is

\begin{equation}
M_{\mathrm{tot}}
=
\sum_{r=1}^{N_B}M_r.
\end{equation}

For a matched single-qubit Pauli observable, the implemented estimator
is

\begin{equation}
\widehat{\langle P_i\rangle}
=
\frac{3}{M_{\mathrm{tot}}}
\sum_{r=1}^{N_B}
\mathbb I(b_{r,i}=P)
\sum_m
(-1)^{z_{r,m,i}},
\label{eq:ibm_shadow_one}
\end{equation}

and for same-axis two-body observables,

\begin{equation}
\widehat{\langle P_iP_j\rangle}
=
\frac{9}{M_{\mathrm{tot}}}
\sum_{r=1}^{N_B}
\mathbb I(b_{r,i}=P)
\mathbb I(b_{r,j}=P)
\sum_m
(-1)^{z_{r,m,i}}
(-1)^{z_{r,m,j}}.
\label{eq:ibm_shadow_two}
\end{equation}

These are the sampled counterparts of the exact PCSR feature set in
Eq.~\eqref{eq:pcsr_pauli_set}.

\begin{breakablealgorithm}
\caption{Random-Pauli Shadow Reconstruction}
\label{alg:ibm_shadow_reconstruct}
\small
\begin{algorithmic}[1]

\Require Measurement sets
$\{\mathbf{Z}^{(r)}\}_{r=1}^{N_B}$,
basis settings
$\{\mathbf{b}^{(r)}\}_{r=1}^{N_B}$,
number of qubits $N$

\Ensure Reconstructed one- and two-body Pauli
expectation estimates

\State Compute total number of shots
\[
M_{\mathrm{tot}}
\gets
\sum_{r=1}^{N_B}
\left|
\mathbf{Z}^{(r)}
\right|
\]

\For{$i=1,\ldots,N$}

    \For{$P\in\{Z,X,Y\}$}

        \State $A\gets0$

        \For{$r=1,\ldots,N_B$}

            \If{$b_i^{(r)}=P$}

                \State
                \[
                A
                \gets
                A+
                \sum_m
                \left(
                1-2Z_{m,i}^{(r)}
                \right)
                \]

            \EndIf

        \EndFor

        \State
        \[
        \widehat{\langle P_i\rangle}
        \gets
        \frac{3A}
        {\max(M_{\mathrm{tot}},1)}
        \]

    \EndFor

\EndFor

\For{$i=1,\ldots,N$}

    \For{$j=i+1,\ldots,N$}

        \For{$P\in\{Z,X,Y\}$}

            \State $A\gets0$

            \For{$r=1,\ldots,N_B$}

                \If{
                $b_i^{(r)}=P$
                \textbf{and}
                $b_j^{(r)}=P$
                }

                    \State
                    \[
                    A
                    \gets
                    A+
                    \sum_m
                    \left(
                    1-2Z_{m,i}^{(r)}
                    \right)
                    \left(
                    1-2Z_{m,j}^{(r)}
                    \right)
                    \]

                \EndIf

            \EndFor

            \State
            \[
            \widehat{\langle P_iP_j\rangle}
            \gets
            \frac{9A}
            {\max(M_{\mathrm{tot}},1)}
            \]

        \EndFor

    \EndFor

\EndFor

\State \Return all reconstructed expectation values

\end{algorithmic}
\end{breakablealgorithm}

\subsection{End-to-end hardware validation}

For each temporal step, $N_B$ random-basis circuits are executed with
$M_B$ shots per basis. The reconstructed feature vector follows the
same ordering as the exact PCSR representation and has dimension
$D_{\mathrm{PCSR}}$ from Eq.~\eqref{eq:pcsr_feature_dim}. Thus,

\begin{equation}
M_{\mathrm{step}}=N_BM_B,
\qquad
M_{\mathrm{all}}=TN_BM_B.
\end{equation}

The hardware-derived features incorporate both finite-shot sampling and
the noise processes represented by the calibrated backend or by the
physical processor.

\begin{breakablealgorithm}
\caption{IBM Heron PCSR Shadow Validation}
\label{alg:run_ibm_pcsr}
\small
\begin{algorithmic}[1]

\Require Topology $\mathcal{T}$,
qubit number $N$,
sequence length $T$,
interaction strength $g$,
number of random bases $N_B$,
shots per basis $M_B$,
random seed $s$

\Ensure Hardware-derived feature matrix
$\mathbf{X}$,
target $\mathbf{y}$,
feature labels $\mathcal{L}$,
total sampler execution time $\tau$

\State
$\mathcal{S}
\gets
\Call{GetIBMSampler}{}$

\State
$\mathcal{R}
\gets
\Call{BuildTopology}{\mathcal{T},N}$

\State Initialize random generator using seed $s$

\For{$i=1,\ldots,N$}

    \State
    $\beta_i^{(z)}
    \sim
    \mathcal{U}(0,2\pi)$

    \State
    $\beta_i^{(x)}
    \sim
    \mathcal{U}(0.3,0.7)$

\EndFor

\State Generate $k$-body Pauli task with $k=2$
\[
(\mathbf{u},\mathbf{y})
\gets
\Call{KBodyPauli}{T,2,s}
\]

\State Set temporal window size
\[
W\gets8
\]

\State Initialize feature collection
$\mathcal{X}\gets[\;]$

\State Initialize total sampler time
$\tau\gets0$

\For{$t=0,\ldots,T-1$}

    \State
    $\ell\gets
    \max(0,t-W+1)$

    \State
    $L_t\gets t-\ell+1$

    \State Initialize zero-padded window
    \[
    \mathbf{w}_t
    \gets
    \mathbf{0}_W
    \]

    \State Right-align
    $\{u_{\ell},\ldots,u_t\}$
    inside $\mathbf{w}_t$

    \State Initialize circuit list
    $\mathcal{C}\gets[\;]$

    \State Initialize basis list
    $\mathcal{B}\gets[\;]$

    \For{$r=1,\ldots,N_B$}

        \For{$i=1,\ldots,N$}

            \State Randomly draw
            \[
            b_i^{(r)}
            \in
            \{Z,X,Y\}
            \]

        \EndFor

        \State
        \[
        C_r
        \gets
        \Call{BuildPCSRCircuit}
        {
        \mathcal{R},
        \mathbf{w}_t,
        g,
        \boldsymbol{\beta}^{(z)},
        \boldsymbol{\beta}^{(x)},
        \mathbf{b}^{(r)}
        }
        \]

        \State Append $C_r$ to $\mathcal{C}$

        \State Append
        $\mathbf{b}^{(r)}$
        to $\mathcal{B}$

    \EndFor

    \State Start execution timer

    \State Submit all circuits simultaneously:
    \[
    J
    \gets
    \mathcal{S}.
    \operatorname{run}
    (\mathcal{C},
    \mathrm{shots}=M_B)
    \]

    \State Wait for sampler results
    $\mathcal{R}_{\mathrm{meas}}$

    \State Update total execution time
    \[
    \tau
    \gets
    \tau+
    \Delta t
    \]

    \For{$r=1,\ldots,N_B$}

        \State
        \[
        \mathbf{Z}^{(r)}
        \gets
        \Call{DecodeBits}
        {
        \mathcal{R}_{\mathrm{meas}}^{(r)},N
        }
        \]

    \EndFor

    \State Reconstruct Pauli estimates
    \[
    \mathcal{E}_t
    \gets
    \Call{ReconstructShadow}
    {
    \{\mathbf{Z}^{(r)}\},
    \{\mathbf{b}^{(r)}\},
    N
    }
    \]

    \State Initialize
    $\mathbf{x}_t\gets[\;]$

    \For{$i=1,\ldots,N$}

        \State Append
        $\widehat{\langle Z_i\rangle}$

        \State Append
        $\widehat{\langle X_i\rangle}$

        \State Append
        $\widehat{\langle Y_i\rangle}$

    \EndFor

    \For{$i=1,\ldots,N$}

        \For{$j=i+1,\ldots,N$}

            \State Append
            $\widehat{\langle Z_iZ_j\rangle}$

            \State Append
            $\widehat{\langle X_iX_j\rangle}$

            \State Append
            $\widehat{\langle Y_iY_j\rangle}$

        \EndFor

    \EndFor

    \State Append
    $\mathbf{x}_t$
    to $\mathcal{X}$

\EndFor

\State Convert feature collection to matrix
\[
\mathbf{X}
\gets
\operatorname{Array}(\mathcal{X})
\]

\State \Return
$\mathcal{L},
\mathbf{X},
\mathbf{y},
\tau$

\end{algorithmic}
\end{breakablealgorithm}
% =====================================================================
\section{Methodological Summary}
\label{sec:method_summary}
% =====================================================================

The methodological hierarchy is now explicit. The hardware-aware
topology and sliding-window encoder define the common reservoir
backbone. The full-CZ model provides the baseline. PCSR generalizes
that baseline by replacing $CZ$ with $CP(2g)$ while leaving the
topology and temporal encoder unchanged. The critical-phase scan then
tests whether the resulting interaction regime is associated with
entanglement, memory, and task performance. Exact and shadow-derived
Pauli features provide the quantum-to-classical interface, LUQPI
provides the supervised readout, classical lag features provide the
non-quantum comparison, and the IBM Heron experiment validates the
same PCSR feature pipeline under sampled noisy execution.

This structure separates definitions from experiments: each mathematical
object is introduced once and later sections refer back to it rather
than redefining it.

\bibliographystyle{elsarticle-num}
% \bibliography{references}

\end{document}
